% ══════════════════════════════════════════════════════════════════
% CONCLUSION GÉNÉRALE ET PERSPECTIVES
% ══════════════════════════════════════════════════════════════════

\chapter*{Conclusion Générale}
\addcontentsline{toc}{chapter}{Conclusion Générale et Perspectives}
\markboth{Conclusion Générale et Perspectives}
{Conclusion Générale et Perspectives}

\section*{Synthèse des travaux}
\addcontentsline{toc}{section}{Synthèse des travaux}

Ce mémoire avait pour ambition de répondre à une question
concrète : comment offrir à l'agriculteur tchadien un outil
d'aide à la décision accessible, fiable et adapté aux réalités
locales ? La réponse apportée est le système \textbf{Agro-IA
	Tchad}, un cadre décisionnel complet articulé autour de quatre
modules.

\medskip

Le \textbf{module Irrigation}, fondé sur un Random Forest
entraîné sur 16~435 observations climatiques réelles de la base
NASA POWER (cinq villes, 2015--2023), prédit le besoin en eau
avec une accuracy de 78,51\,\% et un AUC-ROC de 0,8782,
dépassant le seuil de fiabilité fixé lors de la spécification
des besoins. Le \textbf{module Fertilisation}, s'appuyant sur
XGBoost et 1~599 observations, recommande l'engrais optimal avec
une accuracy remarquable de 95,42\,\%. Le \textbf{module
	Détection des maladies} associe un CNN MobileNetV2 --- entraîné
par transfert d'apprentissage sur un dataset original de 287
images de cultures tchadiennes, une contribution en soi --- à
l'analyse d'indices NDVI issus de Sentinel-2 ou estimés par
caméra smartphone ; son accuracy de 52,27\,\% sur images de
terrain constitue une base solide au regard de la littérature.
Enfin, le \textbf{module Historique} garantit la traçabilité de
toutes les analyses via une base SQLite, avec statistiques,
export CSV et purge sélective, l'ensemble des tests fonctionnels
ayant été validés.

\medskip

L'application web qui intègre ces modules, développée avec
Streamlit et alimentée en temps réel par l'API Open-Meteo, est
déployée sur Streamlit Cloud et accessible gratuitement depuis
tout navigateur. Les objectifs fixés en introduction sont ainsi
atteints.

\section*{Limites rencontrées}
\addcontentsline{toc}{section}{Limites rencontrées}

Ce travail comporte néanmoins des limites qu'il convient de
souligner avec honnêteté. La première tient à la \textbf{taille
	du dataset d'images} : avec environ 28 images par classe, le
module de détection des maladies ne peut rivaliser avec les
systèmes entraînés sur des dizaines de milliers d'images ; les
confusions entre céréales morphologiquement proches (mil,
sorgho) en sont la conséquence directe. La deuxième concerne la
\textbf{persistance sur l'hébergement cloud gratuit}, dont le
système de fichiers éphémère réinitialise la base SQLite à
chaque redémarrage --- l'export CSV atténue cette contrainte
sans la supprimer. La troisième est l'\textbf{absence de
	validation agronomique de terrain} : les recommandations n'ont
pas encore été confrontées à des campagnes réelles menées avec
des agriculteurs et des agronomes de l'ITRAD. Enfin, la
\textbf{dépendance à la connectivité internet} limite l'usage
dans les zones rurales les moins couvertes.

\section*{Perspectives d'amélioration}
\addcontentsline{toc}{section}{Perspectives d'amélioration}

Les perspectives ouvertes par ce travail s'organisent en trois
horizons.

\medskip

\textbf{À court terme}, l'enrichissement du dataset local à
100--200 images par classe --- par des campagnes de collecte
photographique dans les zones de production --- devrait porter
l'accuracy du module maladies à 70--80\,\%. L'intégration de
l'API Sentinel Hub permettrait en outre de récupérer
automatiquement le NDVI d'une parcelle géolocalisée, supprimant
la saisie manuelle.

\medskip

\textbf{À moyen terme}, deux évolutions majeures sont
envisagées : d'une part le passage à une architecture
client-serveur (Django, PostgreSQL) offrant des comptes
utilisateurs et une persistance pérenne de l'historique ;
d'autre part l'intégration de la \textbf{détection des maladies
	par drone multispectral}, qui permettrait de cartographier
l'état sanitaire de parcelles entières à résolution
centimétrique --- capteurs proche infrarouge embarqués, calcul
de cartes NDVI/NDRE complètes et localisation précise des foyers
d'infection --- là où la version actuelle repose sur des
observations ponctuelles (photo de feuille, pixel satellite).

\medskip

\textbf{À long terme}, le système pourrait évoluer vers une
plateforme nationale de suivi agricole : application mobile
fonctionnant hors-ligne avec synchronisation différée, interfaces
en langues locales, intégration de capteurs IoT in situ
(humidité du sol, NPK) et tableaux de bord agrégés à destination
des services agricoles de l'État et de l'ITRAD.

\medskip

En définitive, Agro-IA Tchad démontre qu'une intelligence
artificielle frugale --- données ouvertes, modèles légers, outils
gratuits --- peut produire un outil opérationnel au service de
l'agriculture sahélienne, et pose les fondations techniques d'un
système d'information agricole à l'échelle du pays.